% !TeX program = lualatex
% Bewerbung Masterthesis ABB Forschungszentrum Deutschland – Erlind Sejdiu
% Team: Secure Connected Systems
% Kompilieren mit:  lualatex Anschreiben.tex

\documentclass[a4paper,11pt]{article}

% ─── Pakete ──────────────────────────────────────────────────────────────────
\usepackage{fontspec}
\usepackage[ngerman]{babel}
\usepackage[left=2.2cm, right=2.2cm, top=1.4cm, bottom=1.4cm]{geometry}
\usepackage{microtype}
\usepackage{enumitem}
\usepackage{xcolor}
\usepackage[hidelinks]{hyperref}

% ─── Schrift ─────────────────────────────────────────────────────────────────
\setmainfont{texgyrepagella}[
    Extension = .otf,
    UprightFont = *-regular,
    BoldFont = *-bold,
    ItalicFont = *-italic,
    BoldItalicFont = *-bolditalic
]

\definecolor{ink}{HTML}{1C1C1C}
\definecolor{muted}{HTML}{6E6E6E}
\definecolor{hairline}{HTML}{C4C4C4}
\color{ink}

\hypersetup{pdfauthor={Erlind Sejdiu},
            pdftitle={Bewerbung Masterthesis Secure Connected Systems -- Erlind Sejdiu}}

\setlength{\parindent}{0pt}
\setlength{\parskip}{0.45em}
\pagestyle{empty}

% =============================================================================
\begin{document}

% ─── Kopf (identisch zum Lebenslauf) ─────────────────────────────────────────
{\LARGE Erlind Sejdiu}\\[0.3em]
{\small\color{muted}
Blücherstraße 18, Karlsruhe \enspace\textbar\enspace
0157\,3547\,4746 \enspace\textbar\enspace
\href{mailto:erlindamprogrammieren@gmail.com}{erlindamprogrammieren@gmail.com}}\\[0.4em]
\textcolor{hairline}{\rule{\linewidth}{0.7pt}}

\vspace{0.9em}

% ─── Betreffzeile ────────────────────────────────────────────────────────────
{\large\textbf{Masterarbeit: Benchmark-Konzept \& Tooling für verteilte industrielle Netzwerke}}\\[0.15em]
{\small\color{muted}ABB Forschungszentrum Deutschland \enspace\textbar\enspace Team: Secure Connected Systems \enspace\textbar\enspace Möglicher Beginn: nach Absprache}

\vspace{0.4em}

% ─── Anschreiben ─────────────────────────────────────────────────────────────
Sehr geehrte Damen und Herren,% Falls Ansprechpartner bekannt (z. B. R&D Team Lead): Sehr geehrte(r) Frau/Herr <Name>,

ich studiere Informatik im Master an der Hochschule Karlsruhe mit dem Schwerpunkt Software Engineering und suche eine Masterarbeit, in deren Mittelpunkt die Entwicklung eines Konzepts steht. Ihre Ausschreibung zu einem Benchmark-Konzept und dem zugehörigen Tooling für verteilte industrielle Netzwerke spricht mich daher besonders an, da dies genau meine Stärken anspricht.

Ich habe mich bereits intensiver mit Software-Architekturen auseinandergesetzt. Besonders liegt mir, Anforderungen systematisch zu erfassen, daraus tragfähige Konzepte abzuleiten und diese als Proof of Concept umzusetzen. Das sind genau die Fähigkeiten, auf die es bei der Entwicklung eines Benchmark-Konzepts und des passenden Toolings ankommt.

Wie ich ein solches Konzept methodisch erarbeite, habe ich in meiner Bachelorarbeit bei der ProSeS GmbH (Note 1,0) gezeigt. Dort habe ich ein Konzept zur Verbesserung der Wartbarkeit eines Benachrichtigungssystems entwickelt und umgesetzt. Als Bewertungsgrundlage dienten die Qualitätsmerkmale der ISO/IEC 25010, die ich in messbare Kriterien überführt habe. Dabei habe ich zugleich praktische Erfahrung im industriellen Umfeld gesammelt.

Ergänzend möchte ich auf meine Projektmanagement-Applikation hinweisen, die ich als Microservice-Architektur entwickle. In diesem verteilten System erfasse ich Logs, Traces und Metriken mit OpenTelemetry und werte sie im Aspire-Dashboard aus. So lassen sich das zeitliche Verhalten einzelner Anfragen und das Gesamtverhalten des Systems gezielt analysieren. Erfahrungen, die ich unmittelbar in das Messen und Vergleichen verteilter Systeme einbringen kann.

Statt Sie mit weiterem Text zu überhäufen, überzeuge ich Sie gerne in einem persönlichen Gespräch. Über eine Einladung freue ich mich sehr.

\vspace{0.5em}

Mit freundlichen Grüßen

\vspace{1.2em}

Erlind Sejdiu

\end{document}

